\documentclass{article}
\usepackage{amsthm, amssymb, latexsym, amsmath, amsbsy, graphicx, epsfig}
\usepackage{ mathrsfs }

\title{Writing Assignment \#12: Bijective Functions and Permutations}
\author{Robert T Kohn}
\date{April 2025}

\begin{document}

\maketitle
\section*{Question 1}
\subsection*{Part B}

Claim: $f: (2, \infty) \rightarrow (-\infty, -1)$ given by f(x)=$\frac{-x}{x-2}$ is bijective

One to One:

Suppose $\exists a, b \in (2, \infty) \ni f(a)=f(b)$

Then, $\frac{-a}{a-2}=\frac{-b}{b-2}$

Thus, -a(b-2)=-b(a-2)

So, -ab+2a=-ab+2b

Hence, 2a=2b

Then, a=b

Thus, f is injective

Onto: 

Suppose $\forall y \in (-\infty,-1)$:

Then, y=$\frac{-x}{x-2}$

Thus, x=$\frac{2y}{y+1}$

So, $\exists x \in (2, \infty) \ni x>2$

Hence, $\frac{2y}{y+1}>2$

Then, $2y<2(y+1)$

Thus, $0<2$

So, $(\forall y \in  (-\infty,-1))(\exists x \in (2, \infty)) \ni f(x)=y$

Hence, f is surjective

$\therefore f: (2, \infty) \rightarrow (-\infty, -1)$ given by f(x)=$\frac{-x}{x-2}$ is bijective


\section*{Question 9}
\subsection*{Part B}

$u=\{(1,5),(6,7),(4,4),(5,1),(3,2),(2,6),(7,3)\}$

u=(1 5)(2 6 7 3)


\subsection*{Part F}

$v=\{(1,2),(2,5),(3,4),(5,1),(4,6),(6,7),(7,3)\}$

v(u(1))=v(5)=1

v(u(2))=v(6)=7

v(u(3))=v(2)=5

v(u(4))=v(4)=6

v(u(5))=v(1)=2

v(u(6))=v(7)=3

v(u(7))=v(3)=4

v \textopenbullet u = (2 7 4 6 3 5)


\end{document}
